\section{Results}\label{sec:results}

The results follow the chain. Sections~\ref{sec:res-bank} and \ref{sec:res-persona} ask what the instrument keeps of its anchor: the second-bank check, and a comparison with a country-named prompt and a sign prompt. Sections~\ref{sec:res-screen} to \ref{sec:res-readout} ask how the instrument relates to respondents on new items, in a preregistered screen whose criterion is retrospective. Sections~\ref{sec:res-factorial} and \ref{sec:res-noise} describe the instrument itself. Sections~\ref{sec:res-human} to \ref{sec:res-recovery} report descriptive follow-up analyses that were not preregistered (surface D): agreement with respondents, the distance to their answer distributions, the decomposition of the factorial's cross-talk, and held-out label recovery. Each number carries the label of its surface. Everything computed against the WVS Wave~7 file is retrospective, because that file existed before the instrument.

\subsection{Anchor to instrument: a second scenario bank}\label{sec:res-bank}

The first cohort encodes four of the six declared signs on the strict battery (S1). Over the \dc{42}{countries} countries, the adapter--GPS Spearman correlation is \cl{0.744}{B2_cohort1_trust_rho} \ci{0.500}{0.865}{B2_cohort1_trust_rho} for trust, \cl{0.435}{B2_cohort1_risktaking_rho} \ci{0.085}{0.686}{B2_cohort1_risktaking_rho} for risk taking, \cl{0.488}{B2_cohort1_posrecip_rho} \ci{0.120}{0.695}{B2_cohort1_posrecip_rho} for positive reciprocity, and \cl{0.406}{B2_cohort1_altruism_rho} \ci{0.020}{0.676}{B2_cohort1_altruism_rho} for altruism. These four are the imprint set that the held-out screen later uses.

When the scenario bank is replaced (S2, first column of Table~\ref{tab:checks}), trust keeps its agreement with the GPS, at \cl{0.803}{B2_bank2_trust_rho} \ci{0.615}{0.876}{B2_bank2_trust_rho}. The other three imprint dimensions keep positive point estimates, but their intervals now include zero, so \cl{1}{B2_n_imprint_kept} of the \cl{4}{B2_n_imprint_kept:n} imprint dimensions passes the second-bank check and \cl{3}{B2_n_imprint_dropped} do not. Under the rule fixed before the second cohort was trained, we write no sentence about risk taking, positive reciprocity or altruism carrying over to the second bank. The surfaces that were flat on the first bank stay flat on the second (online supplement).

The two banks agree closely when compared profile by profile. Over the \dc{26}{profiles_realized} shared profiles, the rank correlation between first-bank and second-bank composites is \cl{0.907}{B2_bank_to_bank_trust_spearman} \ci{0.774}{0.957}{B2_bank_to_bank_trust_spearman} for trust. Risk taking is the least stable surface, at \cl{0.544}{B2_bank_to_bank_risktaking_spearman} \ci{0.187}{0.772}{B2_bank_to_bank_risktaking_spearman}, which fits its loss on the second-bank check; the other surfaces agree more closely across banks (online supplement gives both the rank and the Pearson correlation for every surface). Bank-to-bank agreement measures how much the scenario text changes the instrument. It says nothing about people.

\subsection{Anchor to instrument: naming the country or stating the signs}\label{sec:res-persona}

The comparison a survey researcher would reach for first is a prompt that names the population. Scored on the same stack, the country-named persona route does not reproduce the GPS trust order (P): its correlation over the \dc{42}{countries} countries is \cl{0.028}{PX_persona_trust_rho} \ci{-0.339}{0.381}{PX_persona_trust_rho}. The planned contrast with the first cohort is \cl{0.716}{PX_D_cohort1_trust} \ci{0.191}{1.154}{PX_D_cohort1_trust}, so on trust the adapters track the benchmark's order more closely than the named-country prompt, and the reading is the same under a country bootstrap. The sign-profile prompt, which states the same six signs in words, tracks trust as closely as the adapters, at \cl{0.773}{PX_signprompt_trust_rho} \ci{0.584}{0.844}{PX_signprompt_trust_rho}. On trust, what separates the routes is whether the declared sign is supplied, whichever channel carries it. The contrast sets declared signs against a country name and says nothing about trained weights against prompting. On the other surfaces the persona route carries part of the patience ordering, positively on one item and negatively on the other, and its remaining intervals include zero (online supplement).

\subsection{Instrument to respondents: the retrospective screen}\label{sec:res-screen}

Retrospectively, the screen rejects its registered null. The probe selected \cl{56}{SCR_A_primary:n} of the \dc{223}{pool_items} pool items under the primary rule. Their mean signed agreement with the WVS Wave~7 country differences is $A=\cl{0.061}{SCR_A_primary}$, with a country-bootstrap interval of \ci{-0.007}{0.126}{SCR_A_primary}. Under the permutation null, which reassigns the same dimension and sign pairs to items drawn at random from the pool, the null distribution has mean \cl{0.012}{SCR_null_mean} and ninety-fifth percentile \cl{0.046}{SCR_null_q95}, and the one-sided $p$ is \cl{0.010}{SCR_p_perm}. Item--country cells that were never fielded enter as missing (online supplement).

By the registered reading, the signal is weak. Two WVS waves agree with each other on the selected items fielded in both waves with a mean Spearman correlation of \cl{0.746}{SCR_ceiling_mean} (median \cl{0.767}{SCR_ceiling_median}, \cl{34}{SCR_ceiling_mean:n} items). The screen's $A$ is \cl{8.1}{SCR_A_over_ceiling:pct}\% of that ceiling, well below the registered bar of one third, which would be \cl{0.249}{SCR_one_third_ceiling}. The bootstrap interval reaches zero while the permutation test rejects the null because the two answer different questions: whether the assignments beat random assignments of the same pairs to items from the same pool, and how much $A$ varies across samples of countries. The direction of the finding is clearer than its size.

The registered reliability checks point the same way. The secondary rule gives $A=\cl{0.070}{SCR_secondary_A}$ over \cl{185}{SCR_secondary_A:n} items with a one-sided permutation $p\cl{<0.001}{SCR_secondary_p_perm}$, and dropping one country at a time moves the primary $A$ only between \cl{0.051}{SCR_loco_min} and \cl{0.072}{SCR_loco_max}. An exploratory check against each wave's country means, on the items fielded in both waves, shows that the small signal does not depend on one wave (online supplement).

% Table 2: both checks per imprint dimension. Every number is a \cl / \ci macro tied to NUMBERS_TABLE.csv.
\begin{table}[htbp]
\centering
\caption{The second-bank check and the retrospective screen, by imprint dimension. The two checks are reported side by side and never averaged.}\label{tab:checks}
\small
\renewcommand{\arraystretch}{1.2}
\begin{tabularx}{\textwidth}{@{}>{\raggedright\arraybackslash}p{0.17\textwidth}>{\raggedright\arraybackslash}X>{\raggedright\arraybackslash}X>{\raggedright\arraybackslash}X>{\raggedright\arraybackslash}X@{}}
\toprule
 & S2: second bank & H: primary rule & H: primary rule, Bonferroni & H: secondary rule \\
Dimension & $\rho$ [\dc{95}{ci_level}\% interval] & $A_d$ [\dc{95}{ci_level}\% interval], items & $A_d$ [adjusted interval] & $A_d$ [\dc{95}{ci_level}\% interval], items \\
\midrule
Trust & \cl{0.803}{B2_bank2_trust_rho} \newline \ci{0.615}{0.876}{B2_bank2_trust_rho}
 & \cl{0.062}{PD_adapter_primary_trust_A} \newline \ci{-0.047}{0.173}{PD_adapter_primary_trust_A}, \cl{12}{PD_adapter_primary_trust_A:n}
 & \ci{-0.088}{0.199}{PD_adapter_primary_trust_A_bonf}
 & \cl{0.092}{PD_adapter_secondary_trust_A} \newline \ci{-0.019}{0.199}{PD_adapter_secondary_trust_A}, \cl{44}{PD_adapter_secondary_trust_A:n} \\
Risk taking & \cl{0.210}{B2_bank2_risktaking_rho} \newline \ci{-0.199}{0.584}{B2_bank2_risktaking_rho}
 & \cl{0.158}{PD_adapter_primary_risktaking_A} \newline \ci{0.015}{0.290}{PD_adapter_primary_risktaking_A}, \cl{13}{PD_adapter_primary_risktaking_A:n}
 & \ci{-0.026}{0.326}{PD_adapter_primary_risktaking_A_bonf}
 & \cl{0.143}{PD_adapter_secondary_risktaking_A} \newline \ci{0.045}{0.229}{PD_adapter_secondary_risktaking_A}, \cl{50}{PD_adapter_secondary_risktaking_A:n} \\
Positive reciprocity & \cl{0.366}{B2_bank2_posrecip_rho} \newline \ci{-0.068}{0.652}{B2_bank2_posrecip_rho}
 & \cl{-0.071}{PD_adapter_primary_posrecip_A} \newline \ci{-0.281}{0.154}{PD_adapter_primary_posrecip_A}, \cl{11}{PD_adapter_primary_posrecip_A:n}
 & \ci{-0.332}{0.205}{PD_adapter_primary_posrecip_A_bonf}
 & \cl{-0.049}{PD_adapter_secondary_posrecip_A} \newline \ci{-0.198}{0.117}{PD_adapter_secondary_posrecip_A}, \cl{22}{PD_adapter_secondary_posrecip_A:n} \\
Altruism & \cl{0.364}{B2_bank2_altruism_rho} \newline \ci{-0.057}{0.653}{B2_bank2_altruism_rho}
 & \cl{0.069}{PD_adapter_primary_altruism_A} \newline \ci{-0.036}{0.177}{PD_adapter_primary_altruism_A}, \cl{20}{PD_adapter_primary_altruism_A:n}
 & \ci{-0.058}{0.198}{PD_adapter_primary_altruism_A_bonf}
 & \cl{0.040}{PD_adapter_secondary_altruism_A} \newline \ci{-0.038}{0.117}{PD_adapter_secondary_altruism_A}, \cl{69}{PD_adapter_secondary_altruism_A:n} \\
\bottomrule
\end{tabularx}
\par\smallskip
\begin{minipage}{\textwidth}\footnotesize
\textit{Note:} S2 is the Spearman correlation between second-bank adapter composites and the GPS coordinate over \dc{42}{countries} countries, with a profile-block bootstrap (\dc{2,000}{boot_draws} draws, seed \dc{20261011}{seed_c3}); model side only. H is the retrospective screen against WVS Wave~7 country means: $A_d$ is the mean signed item agreement over the selected items the instrument assigned to the dimension, with the registered country bootstrap (\dc{2,000}{boot_draws} resamples, seed \dc{20261015}{seed_screen}). The registered intervals are unadjusted. The Bonferroni intervals over the \dc{4}{imprint_dims} imprint dimensions are exploratory and read off the same draws. The secondary rule is a registered reliability check. No dimension has an interval above zero on both S2 and H.
\end{minipage}
\end{table}


The per-dimension readings show where the signal sits (Table~\ref{tab:checks}). Under the primary rule, only the items assigned to risk taking have an interval above zero, $A_d=\cl{0.158}{PD_adapter_primary_risktaking_A}$ \ci{0.015}{0.290}{PD_adapter_primary_risktaking_A} on \cl{13}{PD_adapter_primary_risktaking_A:n} items, of which \cl{92.3}{PD_adapter_primary_risktaking_share_pos:pct}\% agree in sign. The intervals for trust and altruism include zero. The items assigned to positive reciprocity point the wrong way: their point estimate is negative, and only \cl{18.2}{PD_adapter_primary_posrecip_share_pos:pct}\% of the \cl{11}{PD_adapter_primary_posrecip_share_pos:n} items agree in sign. The secondary rule agrees on both points (online supplement).\footnote{The probe's per-dimension permutation $p$ values were not registered. For altruism under the primary rule they point the other way from the registered interval, so we report the registered intervals only.}

These per-dimension intervals are unadjusted for testing four dimensions. In the exploratory Bonferroni adjustment, the primary-rule interval for risk taking becomes \ci{-0.026}{0.326}{PD_adapter_primary_risktaking_A_bonf} and includes zero, so the one primary-rule dimension that excludes zero does not survive.

The two checks in Table~\ref{tab:checks} disagree by dimension. Trust is the dimension that survives the second bank, and its retrospective screen interval includes zero. Risk taking is the dimension whose screen interval excludes zero before adjustment, and it is one of the three that lose their interval on the second bank. No imprint dimension passes both checks, and the per-dimension cells are small.

\subsection{Instrument to respondents: questionnaire readers and the prompt route}\label{sec:res-comparators}

% Table 3: registered contrasts with the questionnaire comparators and with the prompt route (surface H, retrospective).
\begin{table}[htbp]
\centering
\caption{Retrospective screen: the adapters against two questionnaire readers and against prompting with the same signs.}\label{tab:comparators}
\small
\renewcommand{\arraystretch}{1.2}
\begin{tabularx}{\textwidth}{@{}>{\raggedright\arraybackslash}p{0.24\textwidth}>{\raggedright\arraybackslash}X>{\raggedright\arraybackslash}X>{\raggedright\arraybackslash}X@{}}
\toprule
 & Face-valid reader & Typed decision model, majority of \dc{5}{jev_runs} runs & Sign-profile prompt \\
\midrule
Items assigned by both routes & \cl{5}{CMP_face_valid_D_common:n} & \cl{9}{CMP_jev_D_common:n} & \cl{16}{SP_overlap_with_adapter} \\
Same dimension and sign on those items & \cl{1}{CMP_face_valid_common_same_dim_sign} & \cl{7}{CMP_jev_common_same_dim_sign} & \cl{6}{SP_overlap_same_dim_sign} \\
Adapter minus route, shared items & \cl{0.020}{CMP_face_valid_D_common} \newline \ci{-0.181}{0.215}{CMP_face_valid_D_common} & \cl{-0.002}{CMP_jev_D_common} \newline \ci{-0.064}{0.061}{CMP_jev_D_common} & \cl{0.004}{SP_D_common} \newline \ci{-0.065}{0.075}{SP_D_common} \\
Adapter minus route, own items & \cl{-0.034}{CMP_face_valid_D_own} \newline \ci{-0.242}{0.169}{CMP_face_valid_D_own} & \cl{-0.067}{CMP_jev_D_own} \newline \ci{-0.237}{0.103}{CMP_jev_D_own} & \cl{0.033}{SP_D_own} \newline \ci{-0.083}{0.143}{SP_D_own} \\
Route's own $A$ (items) & \cl{0.095}{CMP_face_valid_A_imprint} (\cl{25}{CMP_face_valid_A_imprint:n}) & \cl{0.128}{CMP_jev_A_imprint} (\cl{40}{CMP_jev_A_imprint:n}) & \cl{0.028}{SP_A_primary} (\cl{56}{SP_A_primary:n}) \\
\bottomrule
\end{tabularx}
\par\smallskip
\begin{minipage}{\textwidth}\footnotesize
\textit{Note:} Differences are adapter minus the named route, with \dc{95}{ci_level}\% intervals from the paired country bootstrap (\dc{2,000}{boot_draws} resamples). The registered primary contrast with each questionnaire reader is the one on shared items. The questionnaire readers' own items are restricted to the imprint set, with the \dc{12}{repaired_items} items repaired by Amendment~1 removed. The adapter's own $A$ is \cl{0.061}{SCR_A_primary} on \cl{56}{SCR_A_primary:n} items. Every interval in the table covers zero.
\end{minipage}
\end{table}


The registered contrasts do not distinguish the adapters from a careful reader of the questionnaire (Table~\ref{tab:comparators}). Every interval covers zero, which the preregistration reads as ``the adapters match a careful reader''. That reading rests on few items. The coverage-matched contrasts use the \cl{5}{CMP_face_valid_D_common:n} items that the adapters and the face-valid rule both assign and the \cl{9}{CMP_jev_D_common:n} items they share with the typed decision model, and the routes agree on dimension and sign for \cl{1}{CMP_face_valid_common_same_dim_sign} and \cl{7}{CMP_jev_common_same_dim_sign} of them. An interval that covers zero on so few items is consistent with a match and does not establish one. Both readers concentrate on trust, where their own-item $A$ exceeds the \cl{0.062}{PD_adapter_primary_trust_A} of the adapters' own trust items, but their per-dimension cells hold too few items to be informative.

Prompting the base model with the same six signs does not pass the registered null. The same probe applied to the \dc{64}{cells} prompted profiles gives $A=\cl{0.028}{SP_A_primary}$ \ci{-0.109}{0.158}{SP_A_primary} with a one-sided permutation $p$ of \cl{0.509}{SP_p_perm}, and it assigns its selected items to a single dimension, trust. The adapter route passes the null and the prompt route does not, but the direct difference between them is unresolved: adapter minus prompt is \cl{0.033}{SP_D_own} on each route's own items and \cl{0.004}{SP_D_common} on the \cl{16}{SP_overlap_with_adapter} items both select, and both intervals cover zero. The two routes agree on dimension and sign for \cl{6}{SP_overlap_same_dim_sign} of those shared items. The paper therefore cannot credit the trained weights over the prompt on the screen.

\subsection{Instrument to respondents: readout robustness}\label{sec:res-readout}

Item-level assignment is stable under option order and paraphrase, whereas membership in the particular selected set churns. The registered statistic, the number of selected items whose dimension and sign are unchanged, is \cl{39}{RO_heldout_reverse_ds_unchanged} of \cl{56}{RO_heldout_reverse_ds_unchanged:n} with reversed options, \cl{39}{RO_heldout_para_a_ds_unchanged} with the first paraphrase and \cl{30}{RO_heldout_para_b_ds_unchanged} with the second. No variant changes more than half, so the registered condition for wording dependence is not met, although the second paraphrase comes close. Membership is another matter: the primary rule is a top-quartile cut on $|t|$, small shifts reorder items near the cut, and only \cl{18}{RO_heldout_reverse_membership_agreement}, \cl{29}{RO_heldout_para_a_membership_agreement} and \cl{16}{RO_heldout_para_b_membership_agreement} of the \cl{56}{RO_heldout_para_b_membership_agreement:n} items are selected under both the original readout and the respective variant (online supplement). A researcher who reads the instrument should trust the dimension and direction it assigns to an item more than the list of items it ranks highest.

\subsection{Inside the instrument: the sign-profile factorial}\label{sec:res-factorial}

% Table 4: the sign-profile factorial on the strict battery (surface F, model side only).
\begin{table}[htbp]
\centering
\caption{Own effects of the six sign bits and summary of entanglement, by route. Strict battery, all \dc{64}{cells} sign profiles.}\label{tab:factorial}
\small
\renewcommand{\arraystretch}{1.15}
\begin{tabularx}{\textwidth}{@{}>{\raggedright\arraybackslash}X r r r r@{}}
\toprule
 & \multicolumn{2}{c}{Trained adapters} & \multicolumn{2}{c}{Sign-profile prompt} \\
\cmidrule(lr){2-3}\cmidrule(l){4-5}
Own effect on surface & Main effect & $t$ & Main effect & $t$ \\
\midrule
Trust & \cl{0.364}{FX_adapter64_trust_own_beta} & \cl{16.74}{FX_adapter64_trust_own_t} & \cl{1.780}{FX_signprompt64_trust_own_beta} & \cl{394.46}{FX_signprompt64_trust_own_t} \\
Negative reciprocity & \cl{0.625}{FX_adapter64_negrecip_own_beta} & \cl{11.64}{FX_adapter64_negrecip_own_t} & \cl{-0.023}{FX_signprompt64_negrecip_own_beta} & \cl{-1.80}{FX_signprompt64_negrecip_own_t} \\
Altruism & \cl{0.242}{FX_adapter64_altruism_own_beta} & \cl{6.32}{FX_adapter64_altruism_own_t} & \cl{0.364}{FX_signprompt64_altruism_own_beta} & \cl{11.99}{FX_signprompt64_altruism_own_t} \\
Patience (Q50) & \cl{0.133}{FX_adapter64_patience_Q50_own_beta} & \cl{6.12}{FX_adapter64_patience_Q50_own_t} & \cl{0.421}{FX_signprompt64_patience_Q50_own_beta} & \cl{17.63}{FX_signprompt64_patience_Q50_own_t} \\
Risk taking & \cl{0.140}{FX_adapter64_risktaking_own_beta} & \cl{4.14}{FX_adapter64_risktaking_own_t} & \cl{-0.457}{FX_signprompt64_risktaking_own_beta} & \cl{-22.31}{FX_signprompt64_risktaking_own_t} \\
Patience (Q43) & \cl{0.049}{FX_adapter64_patience_Q43_own_beta} & \cl{3.63}{FX_adapter64_patience_Q43_own_t} & \cl{-0.069}{FX_signprompt64_patience_Q43_own_beta} & \cl{-19.75}{FX_signprompt64_patience_Q43_own_t} \\
Positive reciprocity & \cl{0.124}{FX_adapter64_posrecip_own_beta} & \cl{3.02}{FX_adapter64_posrecip_own_t} & \cl{0.103}{FX_signprompt64_posrecip_own_beta} & \cl{5.66}{FX_signprompt64_posrecip_own_t} \\
\midrule
Summary of entanglement & \multicolumn{2}{r}{Trained adapters} & \multicolumn{2}{r}{Sign-profile prompt} \\
\midrule
Own effects against the intended direction & \multicolumn{2}{r}{\cl{0}{FX_adapter64_own_effects_wrong_signed} of \cl{7}{FX_adapter64_own_effects_wrong_signed:n}} & \multicolumn{2}{r}{\cl{3}{FX_signprompt64_own_effects_wrong_signed} of \cl{7}{FX_signprompt64_own_effects_wrong_signed:n}} \\
Bits dominant on their own surface (V5) & \multicolumn{2}{r}{\cl{2}{FX_adapter64_V5_hits} of \cl{6}{FX_adapter64_V5_hits:n}} & \multicolumn{2}{r}{\cl{2}{FX_signprompt64_V5_hits} of \cl{6}{FX_signprompt64_V5_hits:n}} \\
Same, on standardized effects & \multicolumn{2}{r}{\cl{2}{FX_adapter64_rowwise_standardized_hits} of \cl{6}{FX_adapter64_rowwise_standardized_hits:n}} & \multicolumn{2}{r}{\cl{1}{FX_signprompt64_rowwise_standardized_hits} of \cl{6}{FX_signprompt64_rowwise_standardized_hits:n}} \\
Surfaces whose largest mover is their own bit & \multicolumn{2}{r}{\cl{1}{FX_adapter64_columnwise_hits} of \cl{7}{FX_adapter64_columnwise_hits:n}} & \multicolumn{2}{r}{\cl{1}{FX_signprompt64_columnwise_hits} of \cl{7}{FX_signprompt64_columnwise_hits:n}} \\
Surfaces whose largest mover is the trust bit & \multicolumn{2}{r}{\cl{3}{FX_adapter64_surfaces_largest_mover_is_trust_bit} of \cl{7}{FX_adapter64_surfaces_largest_mover_is_trust_bit:n}} & \multicolumn{2}{r}{\cl{7}{FX_signprompt64_surfaces_largest_mover_is_trust_bit} of \cl{7}{FX_signprompt64_surfaces_largest_mover_is_trust_bit:n}} \\
Variance on the first principal component & \multicolumn{2}{r}{\cl{51.3}{FX_adapter64_pc1_share:pct}\%} & \multicolumn{2}{r}{\cl{83.5}{FX_signprompt64_pc1_share:pct}\%} \\
Mean absolute correlation between surfaces & \multicolumn{2}{r}{\cl{0.395}{FX_adapter64_mean_abs_inter_surface_r}} & \multicolumn{2}{r}{\cl{0.807}{FX_signprompt64_mean_abs_inter_surface_r}} \\
\bottomrule
\end{tabularx}
\par\smallskip
\begin{minipage}{\textwidth}\footnotesize
\textit{Note:} Main effects are differences in mean surface composite between profiles with the bit set and without it, from an OLS fit with the six bits and their \dc{15}{interactions} two-way interactions (\dc{42}{error_df} error degrees of freedom); composites are expected recoded responses averaged over the surface's items, so the units differ across surfaces. The manipulation is of synthetic training labels or prompted signs; the effects are changes in model response distributions. V5 was a planned acceptance check (at least \dc{4}{v5_required} of \dc{6}{sign_bits} bits); the other summary rows are exploratory.
\end{minipage}
\end{table}


On the trained route, every sign bit moves its own surface in the intended direction (Table~\ref{tab:factorial}). All seven own effects are positive, from \cl{0.364}{FX_adapter64_trust_own_beta} for trust ($t=\cl{16.74}{FX_adapter64_trust_own_t}$) to \cl{0.124}{FX_adapter64_posrecip_own_beta} for positive reciprocity ($t=\cl{3.02}{FX_adapter64_posrecip_own_t}$), and \cl{0}{FX_adapter64_own_effects_wrong_signed} of \cl{7}{FX_adapter64_own_effects_wrong_signed:n} point the wrong way.

The planned separability check fails. Only \cl{2}{FX_adapter64_V5_hits} of the \cl{6}{FX_adapter64_V5_hits:n} bits, those for altruism and negative reciprocity, move their own surface more than any other surface, where at least \dc{4}{v5_required} were required (binomial $p=\cl{0.207}{FX_adapter64_V5_p_binomial}$). The scale-free views agree: \cl{2}{FX_adapter64_rowwise_standardized_hits} of \cl{6}{FX_adapter64_rowwise_standardized_hits:n} on standardized effects, and \cl{1}{FX_adapter64_columnwise_hits} of \cl{7}{FX_adapter64_columnwise_hits:n} surfaces, trust, whose largest mover is its own bit. Each bit moves its own surface and also moves others, often by more: the risk-taking bit moves the altruism surface by \cl{0.983}{FX_adapter64_risktaking_bit_on_altruism}, against the altruism bit's own effect of \cl{0.242}{FX_adapter64_altruism_own_beta}, and the trust bit is the largest mover on \cl{3}{FX_adapter64_surfaces_largest_mover_is_trust_bit} of the \cl{7}{FX_adapter64_surfaces_largest_mover_is_trust_bit:n} surfaces. One common direction carries \cl{51.3}{FX_adapter64_pc1_share:pct}\% of the variance of the seven standardized surfaces, and the mean absolute correlation between surfaces is \cl{0.395}{FX_adapter64_mean_abs_inter_surface_r}. The sign bits are inputs that the instrument keeps distinct in direction and mixes in magnitude.

The prompted route fails the separability check by the same count, \cl{2}{FX_signprompt64_V5_hits} of \cl{6}{FX_signprompt64_V5_hits:n} (binomial $p=\cl{0.207}{FX_signprompt64_V5_p_binomial}$), so V5 does not separate the routes. The common direction and the sign of the own effects do. Under prompting, the first principal component carries \cl{83.5}{FX_signprompt64_pc1_share:pct}\% of the variance, the mean absolute correlation between surfaces is \cl{0.807}{FX_signprompt64_mean_abs_inter_surface_r}, and the trust bit is the largest mover on all \cl{7}{FX_signprompt64_surfaces_largest_mover_is_trust_bit} surfaces. Three prompted own effects point against the intended direction: risk taking at \cl{-0.457}{FX_signprompt64_risktaking_own_beta} ($t=\cl{-22.31}{FX_signprompt64_risktaking_own_t}$), patience Q43 at \cl{-0.069}{FX_signprompt64_patience_Q43_own_beta} ($t=\cl{-19.75}{FX_signprompt64_patience_Q43_own_t}$), and negative reciprocity at \cl{-0.023}{FX_signprompt64_negrecip_own_beta}, which at $t=\cl{-1.80}{FX_signprompt64_negrecip_own_t}$ is not distinguishable from zero. The contrast concerns entanglement and the direction of own effects. It is silent on screening accuracy, where the difference between routes is unresolved.

\subsection{Inside the instrument: retraining noise and the near-zero margin}\label{sec:res-noise}

The unseeded retrain of one profile reproduces the original adapter's readings closely. On the strict battery (\cl{140}{RP_strict_n_cells} scored cells), the largest single-cell difference in code probability is \cl{0.052}{RP_strict_max_abs_prob_diff} and the mean absolute difference is \cl{0.0060}{RP_strict_mean_abs_prob_diff}; the Pearson correlation of expected codes over the \cl{23}{RP_strict_pearson_expected_codes:n} items is \cl{0.9998}{RP_strict_pearson_expected_codes}. On the held-out pool (\cl{1{,}087}{RP_heldout_n_cells} cells), the corresponding values are \cl{0.092}{RP_heldout_max_abs_prob_diff}, \cl{0.0056}{RP_heldout_mean_abs_prob_diff} and \cl{0.9996}{RP_heldout_pearson_expected_codes} over \cl{223}{RP_heldout_pearson_expected_codes:n} items. Every strict-surface shift is at most \cl{0.112}{RP_max_abs_surface_diff_sd} of a standard deviation across profiles, with patience Q50 the largest. This bounds retraining noise for one unit; the variance of retraining across units was not measured.\footnote{Near-zero margin (surface N). The near-zero filter is not promoted under the predeclared rule. Colombia's positive reciprocity moved \cl{0.092}{NZ_COL_posrecip_refit_minus_control_sd:abs} profile standard deviations against its intended direction relative to its control, beyond the declared replicate floor of \cl{0.011}{NZ_COL_posrecip_noise_floor_sd}. A size-matched control, however, differs from the full-data baseline by up to \cl{1.222}{NZ_GBR_max_abs_control_minus_baseline_sd} standard deviations for Great Britain and \cl{1.499}{NZ_COL_max_abs_control_minus_baseline_sd} for Colombia, so the floor is too small to read small-data refits against. We claim no reversal. Two refits are a screen, with no interval and no claim across countries (online supplement).}

\subsection{Instrument to respondents: agreement with the human composite}\label{sec:res-human}

The next four subsections report descriptive follow-up analyses (surface D). They were added after the registered results had been read. None is preregistered, none enters a registered family, and none earns a claim. They set the instrument beside the respondents of the WVS Wave~7 file on the strict battery, and every statement is retrospective. For each surface, Table~\ref{tab:human} gives three Spearman correlations over the \dc{42}{countries} countries: adapter composite with GPS, adapter composite with the human composite of the same items, and the human composite with GPS.

On trust, the first-cohort adapters track the GPS, \cl{0.744}{C9_a_adapter_gps_trust_spearman_rho} \ci{0.500}{0.865}{C9_a_adapter_gps_trust_spearman_rho}, and do not reach the human composite, \cl{0.295}{C9_a_adapter_human_trust_spearman_rho} \ci{-0.044}{0.547}{C9_a_adapter_human_trust_spearman_rho}. All seven adapter--human intervals include zero, and the largest point estimate belongs to negative reciprocity, \cl{0.389}{C9_a_adapter_human_negrecip_spearman_rho} \ci{-0.008}{0.694}{C9_a_adapter_human_negrecip_spearman_rho}. The second-cohort adapters give the same picture on trust, with an adapter--human correlation of \cl{0.340}{C9_b_adapter_human_trust_spearman_rho} \ci{-0.025}{0.586}{C9_b_adapter_human_trust_spearman_rho}; positive reciprocity is the one second-cohort interval that excludes zero, and only just (online supplement).

The link between the GPS and the human composite is itself loose on these items. Its interval excludes zero on four of the seven surfaces (Table~\ref{tab:human}): positive for trust, \cl{0.351}{C9_a_human_gps_trust_spearman_rho} \ci{0.088}{0.564}{C9_a_human_gps_trust_spearman_rho}, and negative for patience Q43, \cl{-0.533}{C9_a_human_gps_patience_q43_spearman_rho} \ci{-0.741}{-0.266}{C9_a_human_gps_patience_q43_spearman_rho}. We read this column as a reference for scale and not as a ceiling, because it does not bound a scorer that is nonlinear or constant within sign profiles. Together, the three columns say that the adapters carry the GPS order on trust, that the GPS order is only loosely related to the human composite, and that the adapters' order of countries does not detectably agree with the human composite on any surface.

% Table: three-column agreement of the first-cohort adapters on the strict battery (surface D; descriptive,
% not preregistered). Every number is a \cl / \ci macro tied to NUMBERS_TABLE.csv.
\begin{table}[htbp]
\centering
\caption{Agreement with the GPS and with the WVS Wave~7 composite, first-cohort adapters, strict battery (surface D, descriptive and not preregistered; retrospective). Spearman correlation over countries, by surface.}\label{tab:human}
\footnotesize
\setlength{\tabcolsep}{4pt}
\renewcommand{\arraystretch}{1.15}
\begin{tabularx}{\textwidth}{@{}>{\raggedright\arraybackslash}X l l l@{}}
\toprule
Surface & Adapter--GPS & Adapter--human & Human--GPS \\
\midrule
Trust & \cl{0.744}{C9_a_adapter_gps_trust_spearman_rho} \ci{0.500}{0.865}{C9_a_adapter_gps_trust_spearman_rho} & \cl{0.295}{C9_a_adapter_human_trust_spearman_rho} \ci{-0.044}{0.547}{C9_a_adapter_human_trust_spearman_rho} & \cl{0.351}{C9_a_human_gps_trust_spearman_rho} \ci{0.088}{0.564}{C9_a_human_gps_trust_spearman_rho} \\
Risk taking & \cl{0.435}{C9_a_adapter_gps_risktaking_spearman_rho} \ci{0.085}{0.686}{C9_a_adapter_gps_risktaking_spearman_rho} & \cl{0.120}{C9_a_adapter_human_risktaking_spearman_rho} \ci{-0.260}{0.426}{C9_a_adapter_human_risktaking_spearman_rho} & \cl{-0.059}{C9_a_human_gps_risktaking_spearman_rho} \ci{-0.429}{0.312}{C9_a_human_gps_risktaking_spearman_rho} \\
Positive reciprocity & \cl{0.488}{C9_a_adapter_gps_posrecip_spearman_rho} \ci{0.120}{0.695}{C9_a_adapter_gps_posrecip_spearman_rho} & \cl{0.324}{C9_a_adapter_human_posrecip_spearman_rho} \ci{-0.074}{0.619}{C9_a_adapter_human_posrecip_spearman_rho} & \cl{0.054}{C9_a_human_gps_posrecip_spearman_rho} \ci{-0.269}{0.386}{C9_a_human_gps_posrecip_spearman_rho} \\
Altruism & \cl{0.406}{C9_a_adapter_gps_altruism_spearman_rho} \ci{0.020}{0.676}{C9_a_adapter_gps_altruism_spearman_rho} & \cl{0.128}{C9_a_adapter_human_altruism_spearman_rho} \ci{-0.170}{0.369}{C9_a_adapter_human_altruism_spearman_rho} & \cl{0.031}{C9_a_human_gps_altruism_spearman_rho} \ci{-0.306}{0.350}{C9_a_human_gps_altruism_spearman_rho} \\
Negative reciprocity & \cl{-0.028}{C9_a_adapter_gps_negrecip_spearman_rho} \ci{-0.418}{0.406}{C9_a_adapter_gps_negrecip_spearman_rho} & \cl{0.389}{C9_a_adapter_human_negrecip_spearman_rho} \ci{-0.008}{0.694}{C9_a_adapter_human_negrecip_spearman_rho} & \cl{0.416}{C9_a_human_gps_negrecip_spearman_rho} \ci{0.173}{0.586}{C9_a_human_gps_negrecip_spearman_rho} \\
Patience (Q43) & \cl{-0.019}{C9_a_adapter_gps_patience_q43_spearman_rho} \ci{-0.414}{0.452}{C9_a_adapter_gps_patience_q43_spearman_rho} & \cl{0.192}{C9_a_adapter_human_patience_q43_spearman_rho} \ci{-0.169}{0.464}{C9_a_adapter_human_patience_q43_spearman_rho} & \cl{-0.533}{C9_a_human_gps_patience_q43_spearman_rho} \ci{-0.741}{-0.266}{C9_a_human_gps_patience_q43_spearman_rho} \\
Patience (Q50) & \cl{0.177}{C9_a_adapter_gps_patience_q50_spearman_rho} \ci{-0.313}{0.580}{C9_a_adapter_gps_patience_q50_spearman_rho} & \cl{0.003}{C9_a_adapter_human_patience_q50_spearman_rho} \ci{-0.291}{0.316}{C9_a_adapter_human_patience_q50_spearman_rho} & \cl{0.388}{C9_a_human_gps_patience_q50_spearman_rho} \ci{0.095}{0.628}{C9_a_human_gps_patience_q50_spearman_rho} \\
\bottomrule
\end{tabularx}
\par\smallskip
\begin{minipage}{\textwidth}\footnotesize
\textit{Note:} Each cell is the Spearman correlation with a \dc{95}{ci_level}\% interval. Over \dc{42}{countries} countries, except the two human columns of risk taking, which cover \cl{41}{C9_a_adapter_human_risktaking_spearman_rho:n} countries because one country has no valid answer to the risk item (Q178). Adapter intervals resample the \dc{26}{profiles_realized} sign profiles (\dc{2,000}{boot_draws} draws, seed \dc{20261011}{seed_c3}); the human--GPS interval resamples countries (seed \dc{20260905}{seed_country}). The human composite is the Wave~7 item-mean composite of the same strict-battery items. The adapter--GPS column repeats the S1 values of Section~\ref{sec:res-bank}.
\end{minipage}
\end{table}


\subsection{Instrument to respondents: distance to the survey distributions}\label{sec:res-tvd}

The correlations above concern the order of countries. A different question is how far the instrument's answer distribution lies from the survey's. For each country--item cell we take the total variation distance between the model's probabilities over the answer codes and the survey-weighted shares of valid human answers, and report its country-equal mean; lower is closer (Section~\ref{sec:design-desc}). Four cells are compared: the base model, the first-cohort country adapters, the country-named persona prompt, and the adapters and the persona prompt together.

The base model's mean distance is \cl{0.470}{C9_TVD_base_mean_tvd_country_equal}. The adapters sit farther from the survey, at \cl{0.515}{C9_TVD_adapter_mean_tvd_country_equal}, and the persona prompt sits closer, at \cl{0.365}{C9_TVD_persona_mean_tvd_country_equal}. Relative to the base model, the adapters are farther by \cl{0.045}{C9_TVD_adapter_minus_base_profile_block} \ci{0.006}{0.093}{C9_TVD_adapter_minus_base_profile_block} under profile blocks, and the persona prompt is closer by \cl{-0.105}{C9_TVD_persona_minus_base_profile_block} \ci{-0.122}{-0.086}{C9_TVD_persona_minus_base_profile_block}; country resampling gives the same readings (online supplement). By dimension, the adapters are farther than the base model on four of the six dimensions and slightly closer on the other two (point estimates; online supplement). With both interventions the distance is \cl{0.420}{C9_TVD_adapter_persona_mean_tvd_country_equal}. The interaction contrast, the adapters' effect with the persona prompt minus their effect without it, is \cl{0.009}{C9_TVD_interaction_profile_block} \ci{-0.003}{0.022}{C9_TVD_interaction_profile_block} under profile blocks, which includes zero, so the data show no detectable departure from additivity; descriptively, the adapters' step away from the survey and the persona prompt's step toward it are consistent with adding.

The ordering of the cells by distance does not follow their ordering by trust agreement with the GPS. The persona prompt has the lowest distance and no trust order (Section~\ref{sec:res-persona}). The adapters have the highest distance and keep the trust order, and adding the persona prompt to the adapters keeps the trust order while the distance falls. The two measures ask different questions, and neither says why the adapters' distributions lie where they do. The instrument carries no information about magnitudes (Section~\ref{sec:discussion}), so a distance to survey distributions is not a quantity it was built to reduce.

\subsection{Inside the instrument: why trust is spared}\label{sec:res-w1}

The factorial shows that each bit moves its own surface in the intended direction and also moves other surfaces (Section~\ref{sec:res-factorial}). To ask why the adapters keep the GPS order of countries on trust and lose it on most other surfaces, we split the main effects of the second-cohort adapters into the effect of each surface's own bit and the summed effect of the other five bits, evaluate each part at the sign profile of each of the \dc{42}{countries} countries, and correlate it with the GPS coordinate of the surface's own dimension (online supplement gives every surface). This is a behavioral decomposition of the adapters' main effects and says nothing about mechanism inside the network.

The own-bit part orders the countries like the GPS on all seven surfaces, with correlations from \cl{0.817}{C9_W1_patience_Q43_spearman_own_vs_gps} to \cl{0.866}{C9_W1_posrecip_spearman_own_vs_gps}. This is mechanical. For a positive own effect, as all seven are, the own part is proportional to the sign of the GPS coordinate, so it restates the correlation between the signs and the GPS and is not a separate test of learning. The full prediction, which adds the cross bits, keeps the GPS order on trust, \cl{0.834}{C9_W1_trust_spearman_full_vs_gps}, and loses most of it elsewhere, with correlations from \cl{-0.020}{C9_W1_patience_Q43_spearman_full_vs_gps} for patience Q43 to \cl{0.422}{C9_W1_altruism_spearman_full_vs_gps} for altruism.

The cross part carries little of the variation of the full prediction on trust and most of it elsewhere. Var(cross)/Var(full) is \cl{0.147}{C9_W1_trust_cross_var_share} on trust and runs from \cl{0.830}{C9_W1_negrecip_cross_var_share} to \cl{1.049}{C9_W1_patience_Q43_cross_var_share} on the other six. This is a ratio of variances and not a percentage of variance explained; it exceeds one on patience Q43 because the own and cross parts covary negatively. Interactions are needed on four surfaces, and they do not restore the GPS order (online supplement gives the covariance shares and the two-way correlations). The rank correlation across the seven surfaces between Var(cross)/Var(full) and the adapters' failure to order countries like the GPS rests on one surface, trust, set against six, so it is not a graded relation and we claim no dose--response (values in the online supplement). No bank label depends on a second dimension: in both banks all \cl{27{,}636}{C9_W1_bank1_label_rows_checked} country-pair training rows follow the sign of the pair's own dimension, with \cl{0}{C9_W1_bank1_label_mismatches} mismatches in the first bank and \cl{0}{C9_W1_bank2_label_mismatches} in the second, so the labels do not carry the cross-talk. We do not claim a cause of the loss on every surface. Component-level analysis, and interventions for any causal statement, have not been done.

\subsection{Anchor to instrument: held-out label recovery}\label{sec:res-recovery}

Every training pair of a country carries a label: the completion that goes with the sign of the country's GPS coordinate on the pair's own dimension is the chosen one. This construction check asks whether that installation carries over to pairs the adapters never saw. For each of the \dc{42}{countries} first-cohort adapters, the margin is the difference between the log probability of the chosen and of the rejected completion. Recovery is the share of held-out pairs on which the adapter's margin exceeds the base model's, accuracy is the share on which the adapter prefers its own chosen completion outright, and base accuracy is the same share for the base model. We report country-equal means with a profile-block bootstrap.

Recovery is \cl{0.940}{C9_label_recovery_recovery_profile_block} \ci{0.926}{0.952}{C9_label_recovery_recovery_profile_block} and accuracy is \cl{0.816}{C9_label_recovery_accuracy_profile_block} \ci{0.771}{0.856}{C9_label_recovery_accuracy_profile_block}. The base model's accuracy is \cl{0.489}{C9_label_recovery_base_accuracy_profile_block} \ci{0.471}{0.506}{C9_label_recovery_base_accuracy_profile_block}, an interval that includes one half, so the base model shows no detectable preference for the chosen completions before training. The mean margin shift is \cl{40.3}{C9_label_recovery_mean_margin_shift_profile_block} \ci{28.2}{54.1}{C9_label_recovery_mean_margin_shift_profile_block} nats. These checks show that training installed each country's synthetic labels on scenarios the adapter never saw. They say nothing about whether the labels order countries like the GPS and nothing about the respondents, and the distance results above show that the same adapters lie farther from the survey's answer distributions than the base model does.
